\documentclass[11pt,a4paper]{article}

\usepackage[margin=2.4cm]{geometry}
\usepackage{amsmath,amssymb,amsfonts}
\usepackage{booktabs}
\usepackage{multirow}
\usepackage{algorithm}
\usepackage{algorithmic}
\usepackage{graphicx}
\usepackage{subcaption}
\usepackage{xcolor}
\usepackage{enumitem}
\usepackage{natbib}
\usepackage{hyperref}
\graphicspath{{../codebase/SkyEngine/experiment/skycausal/figures/}}

\hypersetup{
    colorlinks=true,
    linkcolor=blue!65!black,
    citecolor=blue!65!black,
    urlcolor=blue!65!black
}

\newcommand{\method}{SkyCausal}
\newcommand{\algo}{CBQ-IQL}
\newcommand{\cmax}{C_{\max}}
\newcommand{\cvar}{\operatorname{CVaR}}
\newcommand{\E}{\mathbb{E}}
\newcommand{\Var}{\operatorname{Var}}
\newcommand{\doop}{\operatorname{do}}
\newcommand{\todo}[1]{\textcolor{red}{[TODO: #1]}}

\title{\textbf{SkyCausal: Budgeted Causal Intervention Learning for\\
Rolling-Horizon FJSP--MAPF Scheduling under Path-Execution Uncertainty}}

\author{Hao Wu \and \todo{Co-authors}}
\date{\today}

\begin{document}
\maketitle

\begin{abstract}
Flexible job-shop scheduling with robot transportation requires decisions over
operation sequencing, machine assignment, vehicle assignment, and
collision-free path execution. Existing approaches either replace transportation
with a static travel-time matrix or repeatedly invoke an expensive path planner
without asking when its information can change the production decision. This
paper studies a new \emph{budgeted intervention} problem: given a limited budget
of high-fidelity multi-agent path finding (MAPF) queries and rolling-horizon
rescheduling operations, when should a scheduler intervene, at what granularity,
and how should path-execution uncertainty alter machine and operation decisions?
We propose \method, a closed-loop framework that combines a
domain-constrained dynamic Bayesian causal model with budget-conditioned,
risk-sensitive offline reinforcement learning. The causal model represents how
transport demand and route competition propagate through delivery delay,
machine starvation, and production critical paths. Randomized interventions in
the simulator enable doubly robust estimation of heterogeneous effects for MAPF
queries, AGV reassignment, partial rescheduling, and joint rescheduling.
The resulting posterior effect distributions guide \algo, a conservative
quantile implicit Q-learning policy that minimizes expected and tail makespan
subject to non-discounted computation budgets. A deterministic MAPF safety layer
retains collision-free execution. We define an evaluation protocol covering
causal identification, calibration, compute-quality Pareto efficiency, mediation
analysis, and generalization to unseen maps, fleet sizes, and congestion
mechanisms. \todo{Insert numerical results only after the preregistered
experiments are complete.}
\end{abstract}

\noindent\textbf{Keywords:}
flexible job-shop scheduling; multi-agent path finding; causal inference;
dynamic Bayesian network; offline reinforcement learning; risk-sensitive
scheduling; rolling-horizon optimization

\section{Introduction}
\label{sec:introduction}

Flexible manufacturing systems couple production planning with material
transportation. A flexible job-shop scheduling problem (FJSP) assigns each
operation to a compatible machine and orders operations on shared machines.
When consecutive operations use different machines, an automated guided vehicle
(AGV) or autonomous mobile robot must transport the material through a shared
shop-floor network. The transportation layer is therefore not an independent
post-processing step: machine choices determine where and when traffic is
generated, while route conflicts determine when materials actually reach the
machines.

The literature has progressively moved from fixed transportation times to
joint vehicle scheduling and conflict-aware routing
\citep{berterottiere2024,sun2023,leet2025,li2026,moinuddin2026}.
However, two extremes remain common. A weakly coupled scheduler uses a static
shortest-path matrix and systematically underestimates delay under congestion.
A tightly coupled scheduler repeatedly invokes a high-fidelity MAPF solver and
reoptimizes the production schedule, which may be computationally prohibitive
in a dynamic factory. Neither extreme answers a central operational question:
\emph{when is an expensive path query or rescheduling intervention worth its
cost?}

This question cannot be answered reliably by correlation alone. Historical
logs call the expensive solver more often in difficult states; consequently,
naive regression may associate intervention with poor outcomes even when the
intervention is beneficial. In addition, a locally shorter transport does not
necessarily reduce the final makespan unless the saved time propagates to a
starved machine and then to the production critical path. The relevant quantity
is thus a heterogeneous, system-level causal effect rather than a travel-time
prediction:
\begin{equation}
    \tau_u(s)
    =
    \E\!\left[
        \cmax^{(0)}-\cmax^{(u)}
        \mid S_t=s
    \right],
    \label{eq:cate_intro}
\end{equation}
where \(u\) is a scheduling or routing intervention and \(0\) denotes no
intervention.

Risk is equally important. Vertex conflicts, narrow corridors, and synchronized
task releases create long-tailed transport delays. A method that improves mean
makespan while increasing timeout or deadlock probability is unsuitable for
production. The decision must therefore account for a return distribution, not
only its expectation, while respecting explicit budgets on MAPF calls,
rescheduling calls, and wall-clock time.

We formulate this setting as the \emph{Budgeted Causal FJSP--MAPF
Intervention Problem}. Its distinctive decision is not the low-level motion of
every robot. Instead, a high-level policy chooses among no intervention, a
high-fidelity MAPF query, AGV reassignment, partial production rescheduling, and
joint rolling-horizon rescheduling. Collision-free execution remains the
responsibility of a deterministic planning layer.

We propose \method, with the following intended contributions:

\begin{enumerate}[leftmargin=*]
    \item \textbf{Budgeted causal problem formulation.}
    We formalize selective MAPF evaluation and rolling rescheduling as
    sequential interventions with heterogeneous computational costs. The
    objective jointly controls expected makespan, tail risk, and non-discounted
    query budgets.

    \item \textbf{Cross-layer Bayesian causal model.}
    We introduce a domain-constrained dynamic Bayesian network that represents
    the mechanism from route demand and bottleneck occupancy to delivery delay,
    machine starvation, critical-path extension, and terminal makespan. Its
    posterior separates environmental randomness from epistemic uncertainty and
    supports counterfactual intervention queries.

    \item \textbf{Causal, budget-conditioned distributional policy.}
    We develop Causal Bayesian Quantile Implicit Q-Learning (\algo), which
    combines doubly robust intervention-effect supervision, conservative offline
    policy learning, quantile return estimation, and dynamic budget conditions.
    MAPF feasibility is enforced outside the learned policy.

    \item \textbf{Identification- and mechanism-oriented evaluation.}
    We propose randomized simulator interventions, paired counterfactual
    rollouts, mediation analysis, calibration tests, and mechanism-shift
    generalization, in addition to conventional makespan benchmarks.
\end{enumerate}

These contributions deliberately differ from a generic ``GNN + PPO'' scheduler.
The central hypothesis is that causal propagation and value of intervention
enable better allocation of expensive optimization effort.

\section{Related Work}
\label{sec:related}

\subsection{FJSP with Transportation and Conflict-Free Routing}

FJSP with transportation resources jointly models production operations and
vehicle movements. Berterottiere et al.\ extend disjunctive-graph representations
to transportation resources and introduce dedicated benchmarks
\citep{berterottiere2024}. Sun et al.\ combine a hybrid genetic algorithm with
time-window-based conflict-free routing \citep{sun2023}. Precedence-constrained
task assignment and path finding adds manufacturing task dependencies to MAPF
\citep{liu2024}, while CBS-TA-PTC integrates task assignment, precedence, and
temporal constraints \citep{chong2024}.

Recent work moves closer to joint production and path decisions. ACES jointly
assigns manufacturing processes to machines and generates cyclic robot plans,
with steady-state throughput as the objective \citep{leet2025}. Li et al.\
combine job-shop scheduling and priority-based path planning
\citep{li2026}. Moinuddin et al.\ formulate FJSP with heterogeneous transbots,
zoning, and handoffs using constraint programming \citep{moinuddin2026}.
These works establish the relevance of explicit routing, but they do not study
causal selection of path-planning and rescheduling interventions under a shared
query budget.

\subsection{Bayesian Models for Dynamic Scheduling}

Dynamic Bayesian networks model temporal dependencies and update uncertainty as
new evidence arrives \citep{kungurtsev2024}. Zhu et al.\ use operation- and
system-level Bayesian networks to quantify how disturbances propagate through
an FJSP schedule \citep{zhu2026}. Zheng et al.\ use posterior disturbance impact
to guide dynamic rescheduling under incomplete look-ahead information
\citep{zheng2025}.

Our work differs in three respects. First, uncertainty is generated by
multi-robot path execution rather than only exogenous processing disturbances.
Second, graph nodes include explicit interventions and system-level outcomes.
Third, Bayesian inference is coupled with identifiable treatment-effect
estimation and a long-horizon budgeted policy.

\subsection{Causal Inference and Prescriptive Manufacturing}

Causal machine learning estimates the outcome of interventions rather than
predicting observations. Generalized random forests and doubly robust learners
support heterogeneous effect estimation with flexible nuisance models
\citep{athey2019,chernozhukov2018}. CAnDOIT combines observational and
interventional time-series data for causal discovery \citep{castri2024}.
Recent prescriptive-maintenance work uses intervention distributions and CATE
to rank production-line actions \citep{saretzky2026}. Dynamic targeting under
shared capacity further shows that individual treatment effects must be
adjusted by state-dependent system opportunity costs \citep{hu2025}.

We use the factory simulator as a source of randomized and paired interventions.
This design permits direct validation against simulator ground-truth effects,
which is usually unavailable in observational industrial studies.

\subsection{Distributional, Offline, and Constrained RL for Scheduling}

Most learning-based FJSP methods construct schedules using policy gradients or
graph neural networks. Recent improvement-search methods revise an existing
schedule and can achieve stronger solution quality \citep{wang2025}. Graph-RHO
uses critical-path-aware heterogeneous graphs and adaptive uncertainty
thresholds to accelerate rolling-horizon FJSP \citep{li2026graphrho}. It does not
model explicit multi-robot path execution, identifiable intervention effects,
or shared query budgets. Offline RL reduces simulator interaction and can learn
from solver or heuristic logs. CDQAC uses a conservative discrete quantile
critic for JSP/FJSP and demonstrates learning from suboptimal data
\citep{remmerden2025}.

Distributional RL estimates a return distribution using quantile regression
\citep{dabney2018}; risk-sensitive extensions optimize CVaR or broader spectral
risk measures \citep{moghimi2025}. C2IQL conditions an offline policy on dynamic
constraints while avoiding out-of-distribution action queries
\citep{liu2025c2iql}. O3SRL combines offline RL with online no-regret
optimization for reward-cost trade-offs \citep{chemingui2025}.

\method\ combines these ideas for a new action semantics: deciding whether and
how to invoke costly routing and production interventions.

\section{Problem Formulation}
\label{sec:problem}

\subsection{Factory Model}

Let \(\mathcal{J}\) be a set of jobs, \(\mathcal{M}\) a set of machines, and
\(\mathcal{A}\) a fleet of AGVs operating on graph \(G=(V,E)\). Job \(J_i\)
contains an ordered sequence of operations
\(O_{i1},\ldots,O_{in_i}\). Operation \(O_{ij}\) can be processed by machines in
\(\mathcal{M}_{ij}\), with duration \(p_{ijk}\) on machine \(M_k\).
Machines process at most one operation at a time and processing is
non-preemptive.

If two consecutive operations use different machines, a pickup-delivery task is
released. A path is a time-indexed sequence of vertices. Feasible joint paths
exclude vertex conflicts and opposite-direction edge conflicts and respect
pickup, delivery, and service constraints.

\subsection{Decision Epochs and State}

The simulator advances on an event timeline. A high-level decision epoch occurs
when at least one of the following events happens:

\begin{itemize}[leftmargin=*]
    \item an operation or transport task is released;
    \item a vehicle arrives or becomes idle;
    \item observed transport delay exceeds a prediction interval;
    \item a machine is at risk of inbound-material starvation;
    \item the rolling horizon expires.
\end{itemize}

The pre-intervention state is
\begin{equation}
    S_t =
    \left(
    G_t^{\mathrm{job}},
    G_t^{\mathrm{map}},
    X_t^{\mathrm{machine}},
    X_t^{\mathrm{agv}},
    X_t^{\mathrm{task}},
    H_t,
    B_t
    \right),
    \label{eq:state}
\end{equation}
where \(H_t\) is the logged history and \(B_t=(B_t^q,B_t^r,B_t^c)\) contains
remaining query, rescheduling, and computation budgets. Only variables observed
before intervention are included in \(S_t\).

\subsection{Intervention Actions}

The high-level action is
\begin{equation}
    U_t \in
    \left\{
    u^0,\,
    u^q,\,
    u^a,\,
    u^p,\,
    u^m,\,
    u^j
    \right\},
\end{equation}
representing no intervention, high-fidelity MAPF query, AGV reassignment,
partial sequence repair, machine reassignment, and joint rolling-horizon
rescheduling, respectively. Each action has a vector cost
\(c(U_t)=(c_q,c_r,c_{\mathrm{cpu}})\).

The policy does not output low-level robot moves. Given the selected assignment
and routing request, a deterministic MAPF layer returns paths or an explicit
timeout/infeasibility status.

\subsection{Budgeted Risk-Sensitive Objective}

The primary outcome is terminal makespan including actual path execution:
\begin{equation}
    \cmax = \max_{J_i\in\mathcal{J}} C_i.
\end{equation}

We optimize
\begin{equation}
    \min_\pi
    \E_\pi[\cmax]
    +\lambda_{\alpha}\cvar_{\alpha,\pi}(\cmax)
    +\eta \E_\pi\!\left[\sum_t c_{\mathrm{cpu}}(U_t)\right],
    \label{eq:objective}
\end{equation}
subject to
\begin{equation}
    \sum_t c_q(U_t)\le B^q,\qquad
    \sum_t c_r(U_t)\le B^r,\qquad
    \sum_t c_{\mathrm{cpu}}(U_t)\le B^c.
    \label{eq:budgets}
\end{equation}
Unlike discounted safety proxies, these are episode-level non-discounted
constraints.

\section{The \method\ Framework}
\label{sec:method}

\subsection{Overview}

\method\ contains four modules:

\begin{enumerate}[leftmargin=*]
    \item a closed-loop event simulator that logs production and path execution;
    \item a domain-constrained Bayesian structural causal model;
    \item a randomized intervention and doubly robust effect-estimation layer;
    \item a budget-conditioned distributional offline RL policy with a MAPF
    safety layer.
\end{enumerate}

The simulator first generates observational and interventional trajectories.
The causal model estimates posterior effect distributions. These effects provide
auxiliary supervision and uncertainty penalties for the offline policy. At
deployment, the policy selects only high-level intervention types; conventional
optimization and planning modules execute them.

\subsection{Cross-Layer Dynamic Bayesian Causal Model}

We organize variables into transportation and production layers. Transportation
variables include task-wave intensity \(Q_t\), AGV state \(A_t\), bottleneck
occupancy \(W_t\), candidate-route overlap \(R_t\), transport delay \(D_t\), and
path failure \(F_t\). Production variables include machine queues \(L_t\),
material starvation \(S_t^{m}\), operation releases \(O_t\), critical-path state
\(K_t\), and remaining makespan \(M_t\).

The graph structure combines hard domain constraints and data validation:
\begin{itemize}[leftmargin=*]
    \item job precedence and machine capacity define mandatory directions;
    \item actions cannot have parents from post-intervention variables;
    \item lagged edges are validated using temporal conditional-independence
    tests;
    \item observationally discovered edges that violate process semantics are
    rejected;
    \item randomized interventions validate action-to-outcome directions.
\end{itemize}

A representative factorization is
\begin{align}
    p(X_{t+1}\mid X_t,U_t,\Theta)
    &=
    p(D_{t+1}\mid Q_t,A_t,W_t,R_t,U_t,\Theta_D)
    \nonumber\\
    &\quad
    p(S_{t+1}^{m}\mid D_{t+1},L_t,U_t,\Theta_S)
    \nonumber\\
    &\quad
    p(K_{t+1}\mid S_{t+1}^{m},O_t,\Theta_K)
    \nonumber\\
    &\quad
    p(M_{t+1}\mid K_{t+1},M_t,\Theta_M).
    \label{eq:dbn}
\end{align}

We place priors on \(\Theta\) and maintain a posterior
\(p(\Theta\mid\mathcal{D})\). For nonlinear conditional mechanisms, a neural
conditional density model or deep ensemble may parameterize each factor, while
the graph remains explicit.

\subsection{Randomized and Paired Interventions}

Pure behavior logs are confounded because difficult states trigger more
expensive actions. During data collection, a behavior policy mixes:

\begin{enumerate}[leftmargin=*]
    \item randomized interventions with known propensity;
    \item rule-based policies for operationally plausible states;
    \item current-policy rollouts to cover policy-induced states.
\end{enumerate}

At selected simulator checkpoints, we snapshot the complete state and execute
multiple actions with common random numbers. These paired rollouts approximate
state-level ground-truth potential outcomes while reducing Monte Carlo variance.
They are used for evaluation, not assumed to be available in a real factory.

\subsection{Doubly Robust Heterogeneous Effects}

For clarity, consider binary treatment \(U_t\in\{0,1\}\). Let
\[
    e(s)=P(U_t=1\mid S_t=s)
\]
be the propensity and
\[
    \mu_u(s)=\E[Y_t\mid S_t=s,U_t=u]
\]
the outcome model. The doubly robust pseudo-outcome is
\begin{align}
    \widetilde{\tau}_t
    &=
    \mu_1(S_t)-\mu_0(S_t)
    +\frac{U_t(Y_t-\mu_1(S_t))}{e(S_t)}
    \nonumber\\
    &\quad
    -\frac{(1-U_t)(Y_t-\mu_0(S_t))}{1-e(S_t)}.
    \label{eq:aipw}
\end{align}

Cross-fitting prevents nuisance-model overfitting. Multi-action effects are
estimated relative to \(u^0\), with overlap diagnostics for each action.
The result is a posterior distribution
\begin{equation}
    \tau_u(s)
    \sim
    p\!\left(
    \cmax^{(0)}-\cmax^{(u)}
    \mid S_t=s,\mathcal{D}
    \right).
    \label{eq:effectposterior}
\end{equation}

We use a lower confidence bound
\begin{equation}
    \operatorname{LCB}_u(s)
    =
    \E[\tau_u(s)]
    -\kappa\sqrt{\Var[\tau_u(s)]}
    \label{eq:lcb}
\end{equation}
to avoid interventions whose apparent benefit is supported only in poorly
covered states.

\subsection{Value of a MAPF Query}

A MAPF query is partly informational: it reveals conflict-aware arrival times
and feasibility before committing to a production repair. Let \(Z_t\) denote
the query result. Its expected value of sample information is
\begin{align}
    \operatorname{EVSI}(s)
    &=
    \E_{Z_t\mid s}
    \left[
        \max_{u}
        \E[\mathcal{U}(u)\mid s,Z_t]
    \right]
    \nonumber\\
    &\quad-
    \max_{u}
    \E[\mathcal{U}(u)\mid s]
    -c_{\mathrm{MAPF}}.
    \label{eq:evsi}
\end{align}
The query is useful only when the information can change a downstream
assignment or schedule decision. This separates ``uncertain state'' from
``decision-relevant uncertainty.''

\subsection{\algo}

\algo\ extends implicit Q-learning with three components:

\paragraph{Budget conditioning.}
The policy and value functions receive the remaining budget \(B_t\). Training
randomizes initial budgets, allowing one policy to operate across multiple
compute regimes.

\paragraph{Distributional critic.}
For quantile levels \(\{\rho_i\}_{i=1}^{N}\), the critic predicts
\begin{equation}
    Z_\psi(S_t,U_t,B_t)
    =
    \{q_{\rho_1},\ldots,q_{\rho_N}\}.
\end{equation}
Quantile regression captures tail makespan and permits CVaR evaluation.

\paragraph{Causal and uncertainty guidance.}
The effect posterior supplies an auxiliary target and a pessimism penalty:
\begin{equation}
    \widehat{Q}_{\mathrm{guided}}(s,u,b)
    =
    \widehat{Q}_{\mathrm{offline}}(s,u,b)
    +\lambda_\tau\operatorname{LCB}_u(s)
    -\lambda_{\mathrm{ood}}\sigma_u(s).
    \label{eq:guidedq}
\end{equation}
The causal estimate does not replace the long-horizon critic; it regularizes
action ranking and improves delayed credit assignment.

The total training objective is
\begin{equation}
    \mathcal{L}
    =
    \mathcal{L}_{\mathrm{quantile}}
    +\lambda_V\mathcal{L}_{\mathrm{expectile}}
    +\lambda_{\mathrm{DR}}\mathcal{L}_{\mathrm{effect}}
    +\lambda_{\mathrm{cal}}\mathcal{L}_{\mathrm{calibration}}.
    \label{eq:loss}
\end{equation}

\begin{algorithm}[t]
\caption{\algo\ Training}
\label{alg:training}
\begin{algorithmic}[1]
\REQUIRE Logged trajectories \(\mathcal{D}\), known behavior propensities,
budget distribution \(p(B)\)
\STATE Fit graph-constrained Bayesian mechanisms on training trajectories
\STATE Cross-fit propensity and outcome nuisance models
\STATE Construct doubly robust effect targets for each intervention
\STATE Estimate posterior effect mean, variance, and calibration
\REPEAT
    \STATE Sample transitions and budget conditions from \(\mathcal{D}\)
    \STATE Update quantile critic with distributional Bellman targets
    \STATE Update expectile value without querying out-of-dataset actions
    \STATE Add causal LCB supervision and posterior uncertainty penalty
    \STATE Update budget-conditioned policy by advantage-weighted regression
\UNTIL{validation criteria converge}
\RETURN Policy \(\pi_\phi(U\mid S,B)\), causal model, and calibration report
\end{algorithmic}
\end{algorithm}

\subsection{Safe Execution and Rolling-Horizon Repair}

The intervention policy is not trusted with physical feasibility. A repair
action produces candidate machine, sequence, and vehicle decisions. CP-SAT or
an established FJSP heuristic validates production constraints; PIBT, LaCAM,
PBS, ECBS, or another configured MAPF solver validates paths. Timeouts and
infeasibility are explicit outcomes and consume the configured query budget.

Rolling repair is progressively expressive:
\begin{enumerate}[leftmargin=*]
    \item AGV reassignment only;
    \item sequence repair for unstarted operations;
    \item machine reassignment for unstarted operations;
    \item joint repair within a finite horizon.
\end{enumerate}
This hierarchy permits direct measurement of whether path feedback truly changes
the production decision.

\section{Identification Assumptions}
\label{sec:identification}

We state assumptions rather than treating a predictive graph as automatically
causal.

\paragraph{Consistency.}
Each intervention has a precise implementation, and the observed result under
that implementation equals its potential outcome.

\paragraph{Sequential exchangeability.}
Given logged pre-intervention history \(H_t\), action assignment is independent
of future potential outcomes. Randomized simulator propensities make this
assumption testable for the primary training data.

\paragraph{Positivity.}
Each evaluated intervention has nonzero probability in the states on which its
effect is reported. We report propensity histograms and effective sample size.

\paragraph{Interference.}
Tasks interact through shared machines, paths, and budgets. Therefore, the
treatment unit is the complete system state at an event epoch, not an
independent job or AGV.

\paragraph{Censoring.}
Timeout and unfinished episodes are outcomes. They are never silently removed;
we either model censoring explicitly or assign a preregistered terminal penalty.

\paragraph{Simulator validity.}
Causal conclusions are valid for the simulated structural mechanisms. Transfer
to a physical factory requires calibration or real interventional validation.

\section{Experimental Protocol}
\label{sec:experiments}

\subsection{Research Questions}

\begin{description}[leftmargin=2.0cm,style=nextline]
    \item[RQ1: Coupling]
    When do static transportation times fail, and does explicit path execution
    change machine or sequence decisions?

    \item[RQ2: Causal identification]
    Can the model recover simulator ground-truth intervention effects under
    policy-induced confounding?

    \item[RQ3: Calibration]
    Are posterior transport-delay and effect intervals calibrated on unseen
    instances and maps?

    \item[RQ4: Budgeted control]
    Does \algo\ improve the makespan-runtime Pareto frontier relative to never,
    periodic, threshold, and always-intervene policies?

    \item[RQ5: Tail risk]
    Does distributional control reduce P90/P95 and CVaR makespan without
    excessive mean-performance loss?

    \item[RQ6: Mechanism generalization]
    Does causal structure improve robustness to unseen bottlenecks, fleet sizes,
    task waves, and route solvers?

    \item[RQ7: Mechanism]
    How much benefit is mediated through transport delay, machine starvation,
    and critical-path changes?
\end{description}

\subsection{Instances and Maps}

FJSP instances are drawn from Brandimarte, Hurink, Kacem, Fattahi, and Dauzere
families. Factory layouts include:

\begin{itemize}[leftmargin=*]
    \item open layouts with alternative routes;
    \item single- and double-bottleneck layouts;
    \item warehouse-style intersections;
    \item held-out layouts with changed topology.
\end{itemize}

We vary jobs, machines, AGVs, task-wave synchronization, cross-zone transport,
and path service time. Train, validation, and test splits are separated by both
FJSP instance and map.

\subsection{Routing and Scheduling Backends}

The primary scheduling backends are a dispatching-rule solver and CP-SAT
rolling-horizon solver. Route execution includes:

\begin{itemize}[leftmargin=*]
    \item static BFS travel time;
    \item reactive local A* with deterministic settings;
    \item PIBT or LaCAM as scalable MAPF;
    \item PBS or ECBS as a higher-quality MAPF oracle.
\end{itemize}

All solver calls record status, runtime, conflicts, objective, and timeout.

\subsection{Baselines}

\paragraph{Non-learning policies.}
Never query, always query, periodic query, delay-threshold trigger,
uncertainty-threshold trigger, and myopic expected-value-of-information.

\paragraph{RL policies.}
PPO, QR-DQN or IQN, IQL, constraint-conditioned IQL, and a scheduling-specific
quantile offline RL baseline.

\paragraph{Optimization policies.}
Static FJSP followed by MAPF, transport-time-aware FJSP, fixed rolling horizon,
and small-instance joint CP/MILP where tractable.

\subsection{Ablations}

\begin{itemize}[leftmargin=*]
    \item remove causal graph and use an unconstrained predictor;
    \item replace doubly robust effects with outcome regression;
    \item remove Bayesian posterior and use point estimates;
    \item remove quantile critic and optimize expected return;
    \item replace budget conditioning with a fixed Lagrange multiplier;
    \item remove EVSI and query on uncertainty alone;
    \item allow AGV repair only, without reopening machine assignment;
    \item remove paired counterfactual training/evaluation.
\end{itemize}

\subsection{Metrics}

\paragraph{Production and routing.}
Mean, P90, P95, and CVaR makespan; completion rate; machine starvation; AGV wait
and detour; conflict count; path failure; and deadlock.

\paragraph{Computation.}
MAPF calls, rescheduling calls, total CPU/GPU time, median and tail decision
latency, and timeout rate.

\paragraph{Causal estimation.}
PEHE, ATE/CATE bias, coverage of effect intervals, policy value, overlap,
effective sample size, and calibration error.

\paragraph{Optimality.}
Small-instance optimality gap and budgeted regret relative to a high-budget
joint solver.

\subsection{Statistical Protocol}

Each main condition uses at least ten paired seeds. We report mean, standard
deviation, paired bootstrap confidence intervals, and a preregistered paired
permutation or Wilcoxon test with effect size. Failed and unfinished episodes
remain in the denominator. All methods share solver time limits and either a
fixed wall-clock or oracle-call budget.

\section{Results}
\label{sec:results}

Except for the explicitly labeled routing prevalidation below, this section
remains a preregistered shell. Final numerical claims must be regenerated from
immutable event logs under the complete multi-instance protocol.

\subsection{Does Explicit Path Execution Change the Scheduling Problem?}

\paragraph{Routing-system prevalidation.}
Before training the full causal policy, we tested whether the execution layer
can expose a meaningful compute--quality trade-off. We fixed one
10-job, 50-operation, 6-machine FJSP instance and evaluated 10 maze layouts,
three system seeds, four AGVs, and a 500-step horizon, for 30 paired episodes
per method. Unfinished episodes receive a capped makespan of 500. All methods
use the same greedy production scheduler and seeded random task assigner.

\begin{table}[t]
\centering
\caption{Routing prevalidation over 30 paired episodes. These are system-gate
results, not the final multi-instance evaluation of \algo.}
\label{tab:routing_prevalidation}
\begin{tabular}{lrrrrrr}
\toprule
Method & Success & Job & Operation & Capped \(\cmax\) & Stationary & Query time \\
\midrule
Local A* & 3.3\% & 67.3\% & 81.6\% & 495.8 & 267.9 & 0.00s \\
Rolling LaCAM & 90.0\% & 96.0\% & 98.5\% & 321.4 & 57.0 & 0.27s \\
Adaptive LaCAM & 100.0\% & 100.0\% & 100.0\% & 303.5 & 51.1 & 1.96s \\
\bottomrule
\end{tabular}
\end{table}

Rolling LaCAM replans when task goals or active/idle status changes and
otherwise consumes a cached joint trajectory. Adaptive LaCAM adds a 500ms
LaCAM* refinement only when the initial trajectory exceeds 40 steps; a
candidate replaces the active plan only if it is shorter. Relative to local A*,
the adaptive method reduces capped makespan by 38.79\% and tasked stationary
events by 80.94\%. It triggers 103 refinements, 89 of which shorten the initial
trajectory; five longer candidates are rejected. This establishes the
existence of a useful query-budget trade-off, but does not yet establish
cross-instance or causal-policy superiority.

\begin{figure}[t]
\centering
\begin{subfigure}{0.49\linewidth}
    \includegraphics[width=\linewidth]{routing_completion_rates.pdf}
    \caption{Execution reliability.}
\end{subfigure}
\begin{subfigure}{0.49\linewidth}
    \includegraphics[width=\linewidth]{routing_query_success_pareto.pdf}
    \caption{Query-time versus success.}
\end{subfigure}
\caption{Prevalidation of the routing execution layer.}
\label{fig:routing_prevalidation}
\end{figure}

\begin{figure}[t]
\centering
\includegraphics[width=0.62\linewidth]{routing_paired_makespan_heatmap.pdf}
\caption{Paired capped-makespan reduction of Adaptive LaCAM relative to local
A* across maze and system seeds.}
\label{fig:routing_heatmap}
\end{figure}

\todo{Extend this prevalidation to multiple FJSP families and report the
fraction of instances in which path feedback changes the preferred machine or
sequence decision.}

\paragraph{Preliminary intervention predictor.}
Random counterfactual teacher probes collected 600 refinement outcomes on
states visited by the rolling policy. A PyTorch multi-task MLP predicts
refinement benefit, benefit quantiles, and query cost. On 105 outcomes from two
held-out maps, it obtains 80.0\% benefit classification accuracy, 0.85-step
median-benefit MAE, and 0.047s query-cost MAE. This is predictive
prevalidation only: the network has not yet replaced the threshold online, and
these metrics are not policy-value estimates.

\subsection{Can Intervention Effects Be Identified?}

\todo{Compare naive regression, IPW, doubly robust estimation, causal forest,
and the Bayesian causal model against paired simulator effects.}

\subsection{Main Budgeted Scheduling Results}

\begin{table}[t]
\centering
\caption{Budgeted scheduling performance on held-out instances and maps.}
\label{tab:main}
\begin{tabular}{lccccc}
\toprule
Method & Mean \(\cmax\) & CVaR \(\cmax\) & Completion & MAPF calls & Runtime \\
\midrule
Never query & -- & -- & -- & -- & -- \\
Periodic & -- & -- & -- & -- & -- \\
Always query & -- & -- & -- & -- & -- \\
Threshold & -- & -- & -- & -- & -- \\
C2IQL baseline & -- & -- & -- & -- & -- \\
\algo\ (ours) & -- & -- & -- & -- & -- \\
\bottomrule
\end{tabular}
\end{table}

\todo{Populate Table~\ref{tab:main} from final CSV only.}

\subsection{Causal Mediation and Decision Diagnostics}

\todo{Decompose total effects through transport delay, machine starvation, and
critical-path changes. Show states where a locally beneficial route intervention
does not affect terminal makespan.}

\subsection{Mechanism-Shift Generalization}

\todo{Evaluate held-out bottleneck topology, AGV count, task-wave process, and
route solver. Compare degradation of causal and purely predictive models.}

\section{Implementation Status and Evidence Boundary}
\label{sec:status}

The current SkyEngine prototype contains job, operation, machine, routing-task,
and AGV states; pluggable job, route, and assignment solvers; event-level
transport-to-machine feedback; metrics; and preliminary learning modules. The
initial lifecycle audit exposed a duplicate coordinate offset and local-A*
livelock, with only 8 of 10 jobs and 46 of 50 operations completed. After
repairing the coordinate contract and integrating rolling LaCAM with
collision-semantic alignment, dock reservation, and dock clearance, the
reference lifecycle completes all 10 jobs and 50 operations. The adaptive
prevalidation completes all 30 paired episodes, while the local A*
implementation remains a reactive and incomplete routing baseline.

Consequently, the following are prerequisites for any submitted empirical
claim:

\begin{enumerate}[leftmargin=*]
    \item freeze a versioned multi-instance data contract;
    \item separate deterministic local A* from its stochastic-escape ablation;
    \item validate a second higher-quality MAPF oracle in addition to LaCAM;
    \item mark timeout, fallback, mock, and unfinished episodes explicitly;
    \item regenerate all results under the protocol in
    Section~\ref{sec:experiments}.
\end{enumerate}

This section should be removed from the final submission and replaced by a
reproducibility statement after these gates pass.

\section{Limitations and Threats to Validity}
\label{sec:limitations}

\paragraph{Causal validity.}
Randomized simulator interventions identify effects under the simulator
mechanisms, not automatically in a physical factory. Hidden operational
variables may violate exchangeability in real logs.

\paragraph{Graph misspecification.}
Domain constraints reduce false directions but may omit relevant mechanisms.
We therefore report graph stability and compare against unconstrained outcome
models.

\paragraph{Interference and nonstationarity.}
Shared routes and machines create interference. Policy deployment also changes
the state distribution and potentially the observed causal mechanism. Periodic
recalibration or explicit mechanism-change detection may be required.

\paragraph{Solver dependence.}
The value of querying MAPF depends on solver quality and runtime. Conclusions
must be tested across at least two route solvers.

\paragraph{Scope.}
The first study uses discrete time, unit-speed vehicles, and finite jobs.
Kinematics, battery charging, stochastic arrivals, and physical validation are
future extensions.

\section{Conclusion}
\label{sec:conclusion}

We introduced the Budgeted Causal FJSP--MAPF Intervention Problem and proposed
\method, a framework for deciding when expensive path evaluation and rolling
rescheduling are causally useful. The approach combines an interpretable
cross-layer dynamic Bayesian causal model, randomized and doubly robust
intervention-effect estimation, and a budget-conditioned distributional offline
RL policy. Its intended value is not merely lower predicted travel time, but
better allocation of optimization effort to interventions that propagate to
the production critical path. The proposed evaluation emphasizes causal
identification, posterior calibration, tail risk, compute-quality Pareto
efficiency, and mechanism-shift generalization. Empirical conclusions remain
contingent on the preregistered system and experiment gates.

\begin{thebibliography}{99}

\bibitem[Athey et al.(2019)]{athey2019}
Athey, S., Tibshirani, J., and Wager, S. (2019).
Generalized random forests.
\emph{The Annals of Statistics}, 47(2), 1148--1178.

\bibitem[Berterottiere et al.(2024)]{berterottiere2024}
Berterottiere, L., Dauzere-Peres, S., and Yugma, C. (2024).
Flexible job-shop scheduling with transportation resources.
\emph{European Journal of Operational Research}, 312(3), 890--909.

\bibitem[Castri et al.(2024)]{castri2024}
Castri, L., Mghames, S., Hanheide, M., and Bellotto, N. (2024).
CAnDOIT: Causal discovery with observational and interventional data from time series.
\emph{Advanced Intelligent Systems}, 6(12), 2400181.

\bibitem[Chemingui et al.(2025)]{chemingui2025}
Chemingui, Y., Deshwal, A., Fern, A., Nguyen-Tang, T., and Doppa, J. (2025).
Online optimization for offline safe reinforcement learning.
\emph{Advances in Neural Information Processing Systems}.

\bibitem[Chernozhukov et al.(2018)]{chernozhukov2018}
Chernozhukov, V., Chetverikov, D., Demirer, M., et al. (2018).
Double/debiased machine learning for treatment and structural parameters.
\emph{The Econometrics Journal}, 21(1), C1--C68.

\bibitem[Chong et al.(2024)]{chong2024}
Chong, Y. Q., Li, J., and Sycara, K. (2024).
Optimal task assignment and path planning using conflict-based search with
precedence and temporal constraints.
\emph{Proceedings of AAMAS}.

\bibitem[Dabney et al.(2018)]{dabney2018}
Dabney, W., Rowland, M., Bellemare, M. G., and Munos, R. (2018).
Distributional reinforcement learning with quantile regression.
\emph{Proceedings of AAAI}.

\bibitem[Hu et al.(2025)]{hu2025}
Hu, Y., Li, S., and Wager, S. (2025).
Optimal targeting in dynamic systems.
\emph{arXiv preprint arXiv:2507.00312}.

\bibitem[Kungurtsev et al.(2024)]{kungurtsev2024}
Kungurtsev, V., Rysavy, P., Idlahcen, F., Rytir, P., and Wodecki, A. (2024).
Learning dynamic Bayesian networks from data: Foundations, first principles
and numerical comparisons.
\emph{arXiv preprint arXiv:2406.17585}.

\bibitem[Leet et al.(2025)]{leet2025}
Leet, C., Sciortino, A., and Koenig, S. (2025).
Jointly assigning processes to machines and generating plans for autonomous
mobile robots in a smart factory.
\emph{arXiv preprint arXiv:2502.21101}.

\bibitem[Li et al.(2026)]{li2026}
Li, R., et al. (2026).
A new framework for job shop integrated scheduling and vehicle path planning problem.
\emph{Sensors}, 26(2), 543.

\bibitem[Li et al.(2026)]{li2026graphrho}
Li, Y., Wang, J., Peng, M., Li, G., and Xu, W. (2026).
Graph-RHO: Critical-path-aware heterogeneous graph network for long-horizon
flexible job-shop scheduling.
\emph{arXiv preprint arXiv:2604.10073}.

\bibitem[Liu et al.(2024)]{liu2024}
Liu, S., et al. (2024).
An integrated approach to precedence-constrained multi-agent task assignment
and path finding for mobile robots in smart manufacturing.
\emph{Applied Sciences}, 14(7), 3094.

\bibitem[Liu et al.(2025)]{liu2025c2iql}
Liu, Z., Li, X., and Zhang, J. (2025).
C2IQL: Constraint-conditioned implicit Q-learning for safe offline
reinforcement learning.
\emph{Proceedings of ICML}.

\bibitem[Moghimi and Ku(2025)]{moghimi2025}
Moghimi, M. and Ku, H. (2025).
Beyond CVaR: Leveraging static spectral risk measures for enhanced
decision-making in distributional reinforcement learning.
\emph{Proceedings of ICML}.

\bibitem[Moinuddin et al.(2026)]{moinuddin2026}
Moinuddin, A., Er Raqabi, E. M., and Van Hentenryck, P. (2026).
Scheduling and routing in the flexible job shop with heterogeneous transbots
and zoning: A constraint programming approach.
\emph{arXiv preprint arXiv:2604.24483}.

\bibitem[Remmerden et al.(2025)]{remmerden2025}
van Remmerden, J., Bukhsh, Z., and Zhang, Y. (2025).
Generalizing beyond suboptimality: Offline reinforcement learning learns
effective scheduling through random data.
\emph{arXiv preprint arXiv:2509.10303}.

\bibitem[Saretzky et al.(2026)]{saretzky2026}
Saretzky, F., Andersen, L., Engel, T., and Ansari, F. (2026).
Integrating causal foundation model in prescriptive maintenance framework for
optimizing production line OEE.
\emph{arXiv preprint arXiv:2512.00969}.

\bibitem[Sun et al.(2023)]{sun2023}
Sun, J., et al. (2023).
An approach to integrated scheduling of flexible job-shop considering
conflict-free routing problems.
\emph{Sensors}, 23(9), 4526.

\bibitem[Wang et al.(2025)]{wang2025}
Wang, J., Cao, Z., Zhao, P., et al. (2025).
Learning memory-enhanced improvement heuristics for flexible job shop scheduling.
\emph{Advances in Neural Information Processing Systems}.

\bibitem[Zheng et al.(2025)]{zheng2025}
Zheng, T., Li, D., and Li, J. (2025).
Bayesian dynamic scheduling of multipurpose batch processes under incomplete
look-ahead information.
\emph{arXiv preprint arXiv:2512.01093}.

\bibitem[Zhu et al.(2026)]{zhu2026}
Zhu, Y., Zhang, B., Meng, R., et al. (2026).
Spatiotemporal dual dimensional disturbance perception and Bayesian
quantitative assessment method for flexible job shop scheduling.
\emph{Digital Twin}. DOI: 10.1080/27525783.2026.2690864.

\end{thebibliography}

\end{document}
